%==============================================================================
\section{Act I: The Beginning (\texttt{mono\_euroc.cc})}
%==============================================================================

Our story begins in the \texttt{Examples/Monocular/} folder. This is the \textbf{entry point}---the \texttt{main()} function where everything starts.

\subsection{Scene 1: The Command Line}

When you run ORB-SLAM3 from the terminal:
\begin{lstlisting}[language=bash,numbers=none]
./mono_euroc ../Vocabulary/ORBvoc.txt TUM_mono.yaml /path/to/MH_01 timestamps.txt
\end{lstlisting}

The program expects these arguments:
\begin{enumerate}
    \item \textbf{Vocabulary file}: The ORB dictionary for place recognition
    \item \textbf{Settings file}: Camera calibration, feature parameters, etc.
    \item \textbf{Image folder}: Where the images are stored
    \item \textbf{Timestamps file}: When each image was captured
\end{enumerate}

\subsection{Scene 2: Loading the Images}

The first thing the program does is load the list of images:

\begin{lstlisting}[caption={Loading images in mono\_euroc.cc (simplified)},label=lst:loadimages]
// Load all image paths and timestamps
vector<string> vstrImageFilenames;
vector<double> vTimestamps;

LoadImages("/path/to/MH_01/mav0/cam0/data", 
           "timestamps.txt",
           vstrImageFilenames, 
           vTimestamps);

// Now we have:
// vstrImageFilenames[0] = ".../1403636579763555584.png"
// vstrImageFilenames[1] = ".../1403636579813555456.png"
// ...
\end{lstlisting}

\textbf{Key insight:} At this point, we haven't done any SLAM yet. We're just preparing a list of ``what images to process''.

\subsection{Scene 3: Creating the SLAM System}

This is the most important line in the entire file:

\begin{lstlisting}[caption={Creating the SLAM System}]
ORB_SLAM3::System SLAM(argv[1], argv[2], ORB_SLAM3::System::MONOCULAR, true);
//                     ^vocab    ^config  ^sensor type              ^use viewer
\end{lstlisting}

When this line executes, \textbf{the entire SLAM machinery is created}:
\begin{itemize}
    \item The vocabulary is loaded from disk (takes a few seconds)
    \item Camera calibration is parsed from the YAML file
    \item Three threads are spawned:
    \begin{enumerate}
        \item \textbf{Tracking} (runs in the main thread)
        \item \textbf{LocalMapping} (background thread)
        \item \textbf{LoopClosing} (background thread)
    \end{enumerate}
    \item The Viewer window opens (if enabled)
\end{itemize}

\textit{What happens inside the \texttt{System} constructor? We'll explore this in the next section.}

\subsection{Scene 4: The Main Loop}

After initialization, we enter the main loop---the heartbeat of SLAM:

\begin{lstlisting}[caption={The main processing loop (simplified)}]
for(int ni = 0; ni < nImages; ni++)
{
    // 1. Read the image from disk
    cv::Mat im = cv::imread(vstrImageFilenames[ni], cv::IMREAD_UNCHANGED);
    double timestamp = vTimestamps[ni];
    
    // 2. THE MAGIC LINE: Feed image to SLAM
    SLAM.TrackMonocular(im, timestamp);
    
    // 3. (Optional) Wait to simulate real-time
    // usleep(...);
}
\end{lstlisting}

\textbf{Key insight:} The entire complexity of SLAM is hidden behind one function call: \texttt{TrackMonocular()}. This function takes a raw image and returns the camera pose. Everything else---feature extraction, matching, optimization, map building---happens inside.

\subsection{Scene 5: Cleanup}

After all images are processed:
\begin{lstlisting}[caption={Shutting down}]
SLAM.Shutdown();  // Stop all threads gracefully

// Save trajectory for evaluation
SLAM.SaveTrajectoryEuRoC("CameraTrajectory.txt");
SLAM.SaveKeyFrameTrajectoryEuRoC("KeyFrameTrajectory.txt");
\end{lstlisting}

\subsection{Summary: The Entry Point Flow}

\begin{verbatim}
main()
  |
  +-- Parse command line arguments
  |
  +-- LoadImages() --> Fill vstrImageFilenames, vTimestamps
  |
  +-- Create System SLAM(...) --> Spawns 3 threads
  |
  +-- for each image:
  |       SLAM.TrackMonocular(im, timestamp)
  |
  +-- SLAM.Shutdown()
  |
  +-- SLAM.SaveTrajectory...()
\end{verbatim}
